
\chapter{Outlook}
The research presented in this thesis aimed to characterize the evolution of the Universe through the lens of inverse Optimal Transport (OT). The FLAMINGO simulations were chosen as suitable simulations to work with, and the data was divided into smaller boxes for computational efficiency. The aim was to find a cost matrix that describes the motion of the particles within the simulations. In Chapter~\ref{ch:entropyregiOT}, an algorithm was developed for applying inverse OT, and comments were made on the equivalence of different solutions, as in general inverse OT problems can be ill-posed. This algorithm was inspired by \citet{ma2021}, with their proofs being discretized as can be seen in Appendix \ref{ch:proofs}. The implementation of this algorithm was chosen in such a way to ensure numerical stability and computational efficiency. \par 
Chapter~\ref{ch:results} implements the algorithm, and finds the cost matrix for certain reference cases before embarking on the quest to find the cost matrix for the FLAMINGO data. We looked at linear motion, expansion, collapse and combinations of expansion and collapse. These motions are expected in the universe in voids, clusters, walls and filaments, and therefore we would expect these reference cases to say something about the FLAMINGO cost matrix. All cost matrices were seen to be different for different parameters. The cost matrices between different kinds of motion, for example between expansion and collapse, were seen to have different structures. The cost matrix between the same kind of motion with different parameters often had a different scaling (expansion) or different structure (collapse). This can help in distinguishing the different reference cases in the FLAMINGO cost matrix. \par 
The cost matrix from the FLAMINGO data looks quite symmetric, which is to be expected from the assumption of an isotropic Universe. The hypothesis is that the FLAMINGO cost matrix is a linear combination of multiple reference cases. As a proof of concept to further analyze the cost matrix, a Markov Chain Monte Carlo simulation was made to test this hypothesis with 216 reference cases. However, the Markov chain did not converge, making the resulting cost matrix was quite different from the FLAMINGO cost matrix. The expectation is that more reference cases will solve this issue, but we cannot rule out the possibility that the FLAMINGO data is not a linear combination of the described reference cases or our initial guesses of parameters were too far off.

\section{Future Directions}
First and foremost, a very limited amount of reference cases is considered in this thesis, with the only possibility of expansion and collapse at certain points in the box. More than likely, this is not what is actually going on. For a more representative view, one should have more reference cases, on different scales, or work with another way of interpreting the cost matrix of the data. \par 
Furthermore, a reoccurring theme in this thesis was the computational limits. Hence, future research is suggested to look into the computational efficiency of assigning the data into boxes, or finding the transportation matrix $\pihat$ through an alternative route, achieving better resolution possibly with more particles in larger simulations, possibly recovering the geodesics. Furthermore, future research could focus on different regularization functions $R(c)$, to see if any significantly different, non-equivalent solutions can be found. \par 
The ultimate goal of the framework of inverse OT applied to cosmology is to apply it to larger simulations then we did in this thesis. Then, one could compare the cost matrix of different simulations with different parameters, to eventually compare this theoretical cost matrix to real-world data. Hence, if we compare the real-world data to different simulations with different parameters, the presented framework can be used as a test for General Relativity. However, retrieving any real-world data can be difficult, as to have similar data that was used in the simulations in this thesis, the Cosmic Web must be observed in detail and over a large enough timescale. While clusters of galaxies have been extensively studied and observed, far less is known about filaments, walls and voids \citep{bos2016}. Filaments, walls and voids are by nature hard to observe due to their low density and surface brightness. Hence, observing these objects is difficult, which makes it even more difficult to find the dynamics of these objects. One alternative to trying to observe the full dynamics of the Cosmic Web, might be finding the cost matrix of closer-by, rapidly evolving stellar systems, such as the stars around Sgr A$^*$, the black hole in the galactic center. These stars have been studied extensively, both by numerical simulations and empirical observations, and evolve quite rapidly \citep{gillessen2009, zwart2023}. This makes this system a possible candidate for inverse OT. 


\section{Acknowledgments}
First of all, I would like to thank my Matthieu for introducing me to this project, and his support and idead throughout. I would also like to thank Elena for her interesting suggestions and questions. Besides, I would like to thank Dirk for last-minute stepping up to read through my thesis in the last few weeks of writing. Also, many thanks to Rob McGibbon for helping me with the data reduction when my own Python skills were lacking. Lastly, I want to thank my dad for checking the language of this report. \par 
For making the figures, GitHub Copilot was used to generate the first basic layout of the code.  









